\subsection{From QUBO to Quantum Systems}\label{sec:from-qubo-to-quantum-systems}

So far, we have studied QUBO purely as a mathematical formulation for optimization problems. We have seen how objective functions, Boolean logic, and constraints can all be represented through an energy function whose minimum corresponds to the desired solution. At this point, however, a fundamental question arises: \textbf{\emph{why describe an optimization problem in terms of energy?}}

\highspace
The answer is that \hl{energy is a central concept in physics}. Physical systems naturally evolve toward low-energy configurations, and quantum systems are no exception. If we can encode the energy landscape of a QUBO problem into a quantum system, then \textbf{finding the minimum of the QUBO} becomes \textbf{equivalent} to finding the \textbf{lowest-energy state of that quantum system}.

\highspace
This observation provides the bridge between:
\begin{itemize}
    \item \textbf{Classical Optimization}, represented by a QUBO objective function;
    \item \textbf{Quantum Mechanics}, where the energy of a system is described by a Hamiltonian.
\end{itemize}
The goal of the next sections is therefore to show \hl{how a QUBO energy landscape can be translated into a quantum-mechanical description}. Ultimately, minimizing a QUBO problem will become equivalent to finding the ground state of a suitable Hamiltonian.